\item Find the acceleration \( w \) of body \( 2 \) in the arrangement shown in Fig. 1.15, if its mass is \( \eta \) times as great as the mass of bar \( 1 \) and the angle that the inclined plane forms with the horizontal is equal to \( \alpha \). The masses of the pulleys and the threads, as well as the friction, are assumed to be negligible. Look into possible cases.
    \begin{center}
        \begin{tikzpicture}
            \node at (0, 0) {{image.png}};
        \end{tikzpicture}
    \end{center}
\begin{solution}
    \begin{center}
        \begin{tikzpicture}
            \pic at (0, 0) {frame=3cm};
        \end{tikzpicture}
    \end{center}
    
    \begin{align*}
        \intertext{Let us write Newton's second law for both, bar 1 and body 2 in terms of projection having taken the positive direction of $x_1$ and $x_2$ as shown in the figure and assuming that body 2 starts sliding, say, upward along the incline.}
        T_1 - m_1 g \sin \alpha &= m_1 w_1 & \tag{1}\\
        m_2 g - T_2 &= m_2 w & \tag{2}\\
        \intertext{For the pulley, moving in vertical direction from the equation }
        F_x &= m w_x\\
        2T_2 - T_1 &= (m_p) w_1 = 0 \\
        \intertext{(as mass of the pulley $m_p = 0$)}
        T_1 &= 2T_2 & \tag{3}\\
        \intertext{As the length of the threads are constant, the kinematic relationship of accelerations becomes}
        w &= 2w_1 & \tag{4}\\
        \intertext{Simultaneous solutions of all these equations yields}
        w &= \dfrac{2g \left( 2 \eta - \sin \alpha \right)}{4 \eta + 1} = \dfrac{2g \left( \eta - \sin \alpha \right)}{4 \eta / m_1 + 1}\\
        \intertext{As $\eta > 1$, $w$ is directed vertically downward, and hence in vector form}
        \vec{w} &= \dfrac{2 \left( 2 \eta - \sin \alpha \right)}{4 \eta + 1} \vec{g}
    \end{align*}
\end{solution}